\section{From restricted tiltings to a higher Frobenius functor}
\label{sec:higher}

Fix $r\geq1$, put $q=p^r$, and choose an injective homomorphism
$E=(\mathbb Z/p)^r\hookrightarrow U(k)$.  Write $R$ for restriction
from $SL_2$, and put $\mathcal W=\operatorname{Ver}_{p^r}(k)$.
Theorem~\ref{thm:restricted-tilting} supplies the
representation-theoretic input for the higher Frobenius construction of
Coulembier and Flake.  This section records that construction, explains
how to evaluate it at projectives, and gives the exactness criterion we
will use.

We will repeatedly use the following elementary fact.  In a rigid
abelian tensor category with simple unit, tensoring by a nonzero
object $Y$ is exact and faithful.  Indeed, the coevaluation
$\mathbf1\to Y\otimes Y^*$ is a nonzero monomorphism.  Tensoring
it by $X$ shows that $X\otimes Y=0$ forces $X=0$.  Exactness then
shows that tensoring by $Y$ detects nonzero morphisms through their
images.  An exact faithful functor reflects exactness as well.

\subsection{The functor on group representations}

Let $\mathcal T$ be the category of tilting $SL_2$-modules, and let
$I_r$ be its tensor ideal generated by $\operatorname{St}_r$.
The canonical functor
\[
 \Sigma_r:\mathcal T/I_r\lhook\joinrel\longrightarrow\mathcal W
\]
is fully faithful \cite[Theorem~4.2(i)]{BEO2023}.
For every $s\geq1$, $T(a)$ vanishes modulo $I_s$ exactly when
$a\geq p^s-1$ \cite[Section~3.1]{BEO2023}.
Stable restriction factors through a faithful
functor $\mathcal T/I_r\to\operatorname{Stab}(E)$
\cite[Theorem~5.3.5]{CoulembierFlakeHigher}.
Here $\operatorname{Stab}(E)$ denotes the quotient of
$\operatorname{Rep}_k(E)$ by morphisms factoring through projective
modules.

The assertion that
$R(\operatorname{St}_{r-1})\otimes D$ is a restricted tilting module
for every $D\in\operatorname{Rep}_k(E)$ is precisely the hypothesis
of \cite[Conjecture~5.4.2(1), Theorem~5.4.3]{CoulembierFlakeHigher}.
Since we have proved that assertion, the cited implication produces a
$k$-linear strong symmetric monoidal functor
\begin{equation}\label{eq:theta}
 \Theta:\operatorname{Rep}_k(E)\longrightarrow\mathcal W,
 \qquad \Theta(R(T))\simeq\Sigma_r(T\bmod I_r).
\end{equation}
In particular, $\Theta$ is additive.  It need not be exact.

\begin{lemma}\label{lem:theta-detection}
If $T$ is a tilting module and $R(T)$ is not projective, then
$\Theta(R(T))\ne0$.  Moreover, $\Theta$ annihilates every
$E$-module induced from a proper subgroup of $E$.
\end{lemma}

\begin{proof}
If $T\bmod I_r=0$, stable restriction shows that $R(T)$ is
projective.  Thus nonprojectivity makes $T\bmod I_r$ nonzero;
the fully faithful functor $\Sigma_r$ and \eqref{eq:theta} give the
first assertion.

Stable restriction also shows that $R(\operatorname{St}_r)$ is
projective.  Its dimension is $q$, so it is the regular module:
projectives over the local algebra $kE$ are free.  Equation
\eqref{eq:theta} therefore shows that $\Theta$ kills all projectives.
Set $S=R(\operatorname{St}_{r-1})$.  Its dimension $p^{r-1}$ is
strictly smaller than $q$, so it is nonprojective and $\Theta(S)\ne0$.

Let $H<E$ have rank $s<r$.  For $s>0$, stable restriction at level
$s$ shows that $\operatorname{Res}^E_H S$ is projective, because
$\operatorname{St}_{r-1}\in I_s$.  The assertion is automatic for
$H=1$.  For a $kH$-module $D$, the tensor identity gives
\[
 S\otimes\operatorname{Ind}_H^E D
 \simeq\operatorname{Ind}_H^E
       (\operatorname{Res}^E_H S\otimes D),
\]
which is projective.  Applying $\Theta$ gives
$\Theta(S)\otimes\Theta(\operatorname{Ind}_H^E D)=0$.
Tensoring by a nonzero object of a tensor category is faithful, so
$\Theta(\operatorname{Ind}_H^E D)=0$.
\end{proof}

\subsection{Internal evaluation}

Let $\mathcal C$ be a finite symmetric tensor category.  An object of
$\mathcal C\boxtimes\operatorname{Rep}_k(E)$ is an object of
$\mathcal C$ equipped with an $E$-action.  There is an internal
extension of $\Theta$,
\[
 \Theta_{\mathcal C}:\mathcal C\boxtimes\operatorname{Rep}_k(E)
       \longrightarrow\mathcal C\boxtimes\mathcal W.
\]
We explain its meaning here because $\Theta$ is not assumed exact.
Choose a projective generator $P_0$ and set
$A=\operatorname{End}_{\mathcal C}(P_0)^{\mathrm{op}}$.
Under $\mathcal C\simeq A\text{-}\mathrm{mod}$, an $E$-object is
an $E$-module with a commuting $A$-action.  Applying $\Theta$ to
its underlying $E$-module and to the action maps produces an
$A$-module object in $\mathcal W$, hence an object of
$\mathcal C\boxtimes\mathcal W$.
This is the internal extension used in
\cite[Example~3.2.4, Proposition~3.2.6, and Section~3.3.2]{CoulembierFlakeHigher}; it does not require
right exactness of $\Theta$.

\begin{lemma}\label{lem:projective-evaluation}
For every projective $Q\in\mathcal C$ and every $E$-object $Z$ of
$\mathcal C$, there is a natural isomorphism
\begin{equation}\label{eq:projective-evaluation}
 (\operatorname{Hom}_{\mathcal C}(Q,-)\boxtimes
       \operatorname{Id}_{\mathcal W})(\Theta_{\mathcal C}(Z))
 \simeq\Theta(\operatorname{Hom}_{\mathcal C}(Q,Z)).
\end{equation}
\end{lemma}

\begin{proof}
The $A$-module corresponding to $Q$ is a direct summand of $A^m$
for some $m$.  Thus $\operatorname{Hom}_A(Q,-)$ is obtained from
$m$ copies of the underlying object by applying a split idempotent.
An additive functor preserves finite direct sums and split
idempotents.  Applying this observation to $\Theta$ gives
\eqref{eq:projective-evaluation}, including its naturality.
\end{proof}

\subsection{The tensor-power construction and exactness}

Give $X^{\otimes q}$ the $E$-action arising from regular
permutations of its tensor factors.  Define
\begin{equation}\label{eq:higher-phi}
 \Phi(X)=\Theta_{\mathcal C}(X^{\otimes q})
       \in\mathcal C\boxtimes\mathcal W.
\end{equation}
By Lemma~\ref{lem:theta-detection} and
\cite[Theorem~5.2.1]{CoulembierFlakeHigher}, this construction has
coherent strong symmetric monoidal constraints. They come from the
$E$-equivariant reordering
$(X\otimes Y)^{\otimes q}\simeq X^{\otimes q}\otimes Y^{\otimes q}$,
with the diagonal action on the right, through the internal monoidal
extension of \cite[Section~3.3.2]{CoulembierFlakeHigher}.
The resulting functor is additive and satisfies
\[
 \Phi(af)=a^q\Phi(f)\qquad(a\in k).
\]
Additivity can also be seen directly from the regular action: in the expansion of
$(X\oplus Y)^{\otimes q}$, every nonconstant coloring of the
tensor positions has a proper stabilizer in $E$.  The sum over
its orbit is induced from that stabilizer, so is annihilated.
The two constant colorings give $\Phi(X)$ and $\Phi(Y)$.
The corresponding expansion on morphisms proves additivity,
and scalar multiplication gives the displayed semilinearity.

Combining \eqref{eq:higher-phi} and
Lemma~\ref{lem:projective-evaluation} yields
\begin{equation}\label{eq:phi-evaluation}
 (\operatorname{Hom}_{\mathcal C}(Q,-)\boxtimes
      \operatorname{Id}_{\mathcal W})(\Phi(X))
 \simeq\Theta(\operatorname{Hom}_{\mathcal C}
                       (Q,X^{\otimes q})).
\end{equation}

\begin{proposition}\label{prop:higher-exactness}
If there are projective objects $P,Q\in\mathcal C$ for which
\[
 \Theta(\operatorname{Hom}_{\mathcal C}(Q,P^{\otimes q}))\ne0,
\]
then $\Phi$ is exact.
\end{proposition}

\begin{proof}
Equation~\eqref{eq:phi-evaluation} implies $\Phi(P)\ne0$.
Tensoring a short exact sequence by $P$ makes it split, since
tensor products of projectives with arbitrary objects are
projective.  The additive functor $\Phi$ preserves this split
sequence.  Its monoidal structure identifies the result with
the image sequence tensored by $\Phi(P)$.
Tensoring by the nonzero object $\Phi(P)$ is exact and faithful,
and therefore reflects exactness.  Thus the image sequence is
exact. This argument is the nonzero-projective-image criterion of
\cite[Proposition~7.6 and Remark~7.7(i)]{EtingofOstrikFrobenius},
applied to the higher construction as in
\cite[Proposition~5.2.4]{CoulembierFlakeHigher}.
\end{proof}

The next step is now concrete: we will choose $r$ and the
embedding of $E$ so that
\[
 \Theta(\operatorname{Hom}_{\mathcal C}(P,P^{\otimes q}))\ne0
\]
for a projective generator $P$.
